\section{Models, Protocol and Checks}
\label{sec:setup}

\paragraph{Models, from simplest to richest.}
We organise our own models in stages. \textbf{Stage 0} holds the classical
forecasts: persistence ($\hat v=v_{\mathrm{pre}}$), the training-set mean of the
target (a constant), EWMA ($\lambda{=}0.94$), HAR-RV \citep{corsi2009har} and
GARCH(1,1). \textbf{Stage 1} is a LightGBM model on past-volatility features,
transcript length, number of Q\&A turns and a \textbf{company fixed effect}, our
identity-only baseline. \textbf{Stage 2} uses the transcript: BGE-M3
\citep{bge2024} embeddings pooled per section (1024-d), FinBERT \citep{finbert}
sentiment aggregates and surface statistics. \textbf{Stage 3} uses the audio, resampled to
16\,kHz mono: eGeMAPSv02 functionals \citep{opensmile}, and WavLM-Large
\citep{wavlm} and emotion2vec+ \citep{emotion2vec} embeddings over 30-second
windows, averaged per call. \textbf{Stage 4} combines modalities by gated fusion or by
stacking (a ridge over the text, audio and Stage-1 predictions, fit on the
validation set). From Stage~2 on, every model is reported with and without a block
of past-volatility features ($v_{\mathrm{pre}}$ at all horizons plus the HAR-RV
inputs). The prediction heads are a ridge regression, with its penalty chosen on
validation, and a small MLP run with five seeds. The MLPs are unstable across seeds
(standard deviation of \Rtwo{} up to 6.8), so the main grid uses one deterministic
model per stage, chosen before the results were known; MLP cells are in
Appendix~\ref{app:complete}.

\paragraph{Metrics and tests.}
We report MSE on log realised volatility and, as the headline metric, \Rtwo{}
relative to persistence \citep{gukellyxiu2020}, defined as
$1-\sum(y-\hat y)^2/\sum(y-\hat y_{\mathrm{persist}})^2$. Because $\Delta v$ has a
positive mean on the measured anchor, persistence (which predicts $\Delta v{=}0$)
is a weak reference there, so every $\Delta v$ claim is also compared with the
constant model. Diebold--Mariano tests \citep{dm1995} with the small-sample
correction compare per-call squared errors against persistence and against the
previous stage. The confirmatory comparisons (each stage against Stage~1 on
$\Delta v$ across horizons, Stage~4 against the best single modality, and the
published split against the controlled split for each reproduced model) are
Holm-corrected; all other analyses are exploratory.

\paragraph{Identity checks.}
\label{sec:controls}
We use four. (1)~A \emph{company-only} model predicts each company's training
mean. (2)~A \emph{same-company swap} gives each test call the transcript or audio
of another call by the same company, and a \emph{random swap} gives it any other
call. A model that does not change under the first is reading identity; one that
does not change under the second is reading nothing. (3)~A linear \emph{probe}
tries to recover the ticker from the frozen call embedding. (4)~Every table shows
\tickerdisjoint{} results next to temporal ones. The command
\texttt{ecvol audit predictions} runs the prediction-side versions of these
checks, plus split integrity, seen-versus-unseen companies, the between-company
share of prediction variance and company-clustered significance, on any file with
one row per (call, ticker, date, split, horizon, truth, prediction), so a published
result can be checked from the one file its authors already have.

\paragraph{Reproductions and post-cutoff tests.}
Of seventeen papers in this line, one releases code that runs as released
(Same-Company-Same-Signal), seven release code that needs repairs or inputs that
were never shipped, and nine release nothing usable (Table~\ref{tab:code-audit}).
We rerun the five with usable code (HTML, KeFVP, Same-Company-Same-Signal,
\citealp{sawhney2020mm}, DialogueGAT) \emph{on their own data}, using the authors'
code wherever it can be made to start and a line-by-line port where it cannot
(every patch is in a released ledger), first against the published number
and then under a gapped split and a \tickerdisjoint{} split of the same data.
Earnings25 allows two more tests: the frozen Stage-3 audio models rerun by
bitrate, and the whole FinCall-trained pipeline scored, with no retraining, on
calls that postdate every model's training cutoff.
